% Ad Familiares 15.11 --- headnote
% ad-familiares-15-11, August 50 BC, Tarsus

\textit{Cicero (still in his Cilician province, writing from Tarsus) to
C. Claudius Marcellus, consul of 50 BC. The Perseus dateline assigns the
letter to ``the same place and time'' as \textit{Fam.} 15.6 ---
i.e. roughly August 50, in the closing weeks of Cicero's governorship.
Marcellus had taken a leading role in pressing the Senate to decree
a \textit{supplicatio} (a public thanksgiving) for Cicero's modest
military operations on Mount Amanus, against considerable resistance
from those who thought the honour disproportionate to the campaign.
This short note is his thank-you. The tone is warm but stylized: the
debt Cicero owes the Claudii Marcelli as a family is already an old
one, and what Marcellus has just done as consul renews and intensifies
it.}

\textit{The letter is also a quiet pre-announcement. Cicero closes by
hoping that the voyage home --- now running squarely into the Etesian
winds that blow against any westbound sail in late summer --- will not
hold him too long, and that he will see Marcellus ``before long.''
The text in section 1 is corrupt at one point (Mendelssohn and
Shackleton Bailey both obelize the phrase that closes the section);
the translation supplies a plausible filling in brackets.}
